% Written By Enoch Yu
% Downloaded from archive.enochyu.com

\documentclass{article}

\usepackage{amsmath}
\usepackage{amssymb}
\usepackage{amsthm}
\usepackage{mathtools}

\usepackage[margin=1in, headsep=15pt]{geometry}
\usepackage{lastpage}
\usepackage{fancyhdr}
\fancyhead[R]{Enoch Yu}
\fancyfoot[C]{Page \thepage \hspace{1pt} of \pageref{LastPage}}
\pagestyle{fancy}
\usepackage{parskip}

\usepackage{enumitem}
\usepackage[dvipsnames]{xcolor}
\usepackage[normalem]{ulem}
\usepackage{url}
\usepackage{hyperref}
\hypersetup{
    colorlinks=true,
    linkcolor=black,
    filecolor=magenta,      
    urlcolor=cyan,
    pdftitle={7.8.2025 AMC 10, 12 Free Class Preview Homework Solution},
    pdfpagemode=FullScreen,
}
\usepackage{tabularx}
\usepackage{arydshln}
\usepackage{float}
\usepackage{pdfpages}

\usepackage{tikz}
\usetikzlibrary{calc, angles, shapes.geometric}
\usepackage{tkz-euclide}
\usepackage{pgfplots}
\pgfplotsset{compat=1.9}
\usepackage[font=small,labelfont=bf]{caption}

\newtheorem{theorem}{Theorem}[section]
\newtheorem{lemma}[theorem]{Lemma}
\newtheorem*{lemma*}{Lemma}
\newtheorem{sublemma}{Lemma}[section]
\newtheorem{proposition}{Proposition}
\newtheorem{corollary}{Corollary}[theorem]
\newtheorem{example}{Example}[section]
\newtheorem*{example*}{Example}
\newenvironment{solution}{\begin{trivlist}\item[]{\bf Solution}}{\qed \end{trivlist}}


\title{7.8.2025 AMC 10, 12 Free Class Preview Homework Solution}
\author{Enoch Yu}
\date{July 2025}

\begin{document}

\maketitle

\subsection*{Problem 1.}
Solve $\sqrt{2x + 1} + 2\sqrt{x - 1} = 4$.

\begin{solution}
Let $a = \sqrt{2x + 1}$ and $b = \sqrt{x - 1}$ where
$b \le 2$. Therefore, $a^2 - 2b^2 = 3$ and $a + 2b = 4$.
Substitution could be utilized.
\begin{align*}
    a &= 4 - 2b \\
    (4 - 2b)^2 - 2b^2 &= 3 \\
    2b^2 - 16b + 13 &= 0 \\
    \therefore b &= \frac{8 - \sqrt{38}}{2}
\end{align*}
In other words, $\sqrt{x - 1} = \frac{8 - \sqrt{38}}{2}$.
Meaning, $x = \frac{102 - 16\sqrt{38}}{4} + 1$,
or $\frac{53 - 8\sqrt{38}}{2}$.
\begin{lemma*}
There exists at most one real root to the equation.
\end{lemma*}
\begin{proof}
Let $f(x) = \sqrt{2x + 1} + 2\sqrt{x - 1} - 4$.
\begin{align*}
    f'(x)
    &= \frac{2}{2\sqrt{2x + 1}} + \frac{2}{2\sqrt{x - 1}} \\
    &= \frac{1}{\sqrt{2x + 1}} + \frac{1}{\sqrt{x - 1}}
\end{align*}
Since $f'(x)$ is always positive wherever it is defined,
function $f(x)$ is always increasing. In other words,
there could be at most one real root to the equation.
\end{proof}
\noindent
Therefore, $x = \frac{53 - 8\sqrt{38}}{2}$ is the only real roots.
\end{solution}

\subsection*{Problem 2.}
Solve $\sqrt{x - \frac{5}{x}} - \sqrt{5 - \frac{5}{x}} = x$.
\begin{solution}

\textbf{Key Word} Average Method

Notice that every number could be written as a sum of
two numbers that have same distance from the half the
value of the number. Therefore, let $\sqrt{x - \frac{5}{x}}
= \frac{x}{2} + k$ and $-\sqrt{5 - \frac{5}{x}} = \frac{x}{2} - k$.
\begin{align*}
    \sqrt{x - \frac{5}{x}} &= \frac{x}{2} + k \\
    -\sqrt{5 - \frac{5}{x}} &= \frac{x}{2} - k \\
    x - \frac{5}{x} &= \frac{x^2}{4} + kx + k^2 \\
    5 - \frac{5}{x} &= \frac{x^2}{4} - kx + k^2 \\
    x - 5 &= 2kx \\
    \therefore k &= \frac{x - 5}{2x}
\end{align*}
Substitution could be used.
\begin{align*}
    \sqrt{x - \frac{5}{x}} &= \frac{x}{2} + \frac{x - 5}{2x} \\
    2\sqrt{x - \frac{5}{x}} &= x + \frac{x - 5}{x} \\
    2\sqrt{x - \frac{5}{x}} &= 1 + \left( x - \frac{5}{x} \right) \\
    1 + - 2\sqrt{x - \frac{5}{x}}
        + \left( x - \frac{5}{x} \right) &= 0 \\
    (1 - \sqrt{x - \frac{5}{x}})^2 &= 0 \\
    x - \frac{5}{x} &= 1 \\
    \therefore x &= \boxed{\frac{1 - \sqrt{21}}{2}}
\end{align*}
Through graphing, it is evident that the $x$ value is
the only possible $x$.
\end{solution}

\subsection*{Problem 3.}
Find $x^6 - 2\sqrt{2}x^5 - 3x^4 - x^3 + 2 \sqrt{5}x^2
- 4x + \sqrt{5}$ if $x = \sqrt{5} + \sqrt{2}$.
\begin{solution}
\begin{align*}
    &\quad\ \, x^6 - 2\sqrt{2}x^5 - 3x^4 - x^3
        + 2\sqrt{5}x^2 - 4x + \sqrt{5}\\
    &=x^5 \left( x - 2\sqrt{2} \right) - 3x^4
        + x^2 \left( -x + 2\sqrt{5} \right) - 4x + \sqrt{5} \\
    &= \left( \sqrt{5} - \sqrt{2} \right)(x^5 + x^2)
        - 3x^4 - 4x + \sqrt{5} \\
    &= x\left( \sqrt{5} - \sqrt{2} \right)(x^4 + x)
        - 3(x^4 + x) - x + \sqrt{5} \\
    &= \left( \sqrt{5} + \sqrt{2} \right)
            \left( \sqrt{5} - \sqrt{2} \right)
            - 3(x^4 + x) - \sqrt{2} \\
    &= \boxed{\sqrt{2}}
\end{align*}
\end{solution}

\subsection*{Problem 4.}
Show that $\sqrt[3]{ax^2 + by^2 + cz^2} =
\sqrt[3]{a} + \sqrt[3]{b} + \sqrt[3]{c}$ if
$ax^3 = by^3 = cz^3$ and $\frac{1}{x} + \frac{1}{y}
+ \frac{1}{z} = 1$.
\begin{proof}
\begin{align*}
    \sqrt[3]{ax^2 + by^2 + cz^2}
    &= \sqrt[3]{\frac{ax^3}{x} + \frac{by^3}{y} + \frac{cz^3}{z}} \\
    &= \sqrt[3]{ax^3} \\
    &= \sqrt[3]{ax^3}
        \left( \frac{1}{x} + \frac{1}{y} + \frac{1}{z} \right) \\
    &= \sqrt[3]{a} + \sqrt[3]{b} + \sqrt[3]{c}
\end{align*}
\end{proof}
\end{document}

